\cleardoublepage
\renewcommand{\bibname}{Literatura}
\phantomsection
\addcontentsline{toc}{chapter}{Literatura}
\begin{thebibliography}{99}
\small\raggedright

\bibitem{pac2024} M.~Jasiukowicz, M.~Sadowska, J.~Radgowski, M.~Gwozdowski, M.~Kapciński, A.~Kaczyński, N.~Netkowski, D.~Bagińska, W.~Sieklicki, \emph{Life after BEXUS: Stratospheric balloon activities at Gdańsk University of Technology}, 26th ESA Symposium on European Rocket and Balloon Programmes and Related Research (PAC), Lucerna 2024.\newline\url{https://yasiu.pl/blog/life-after-bexus-pac-2024-in-lucerne/}

\bibitem{iac2023} M.~Jasiukowicz, M.~Sadowska, M.~Grzybowski, B.~Rybacki, K.~Pelzner, S.~Magrian, K.~Loret, W.~Sieklicki, \emph{SimLE Stardust: How an experiment evolved from student tinkering to a stratospheric research platform}, 74th International Astronautical Congress, Baku 2023, IAC-23,A1,6,4,x80448.\newline\url{https://yasiu.pl/blog/seven-years-of-stardust-iac-2023-in-baku/}

\bibitem{raport2020} Zespół Stardust, \emph{Test Mission Technical Report 2020/07/20}, SimLE 2020.\newline\url{https://github.com/simle-stardust/Manual/blob/main/papers/2020_Stardust_Test_Mission_Report.pdf}

\bibitem{praca2022} M.~Jasiukowicz, \emph{Stardust -- Stratosferyczna platforma eksperymentalna}, praca dyplomowa inżynierska, Politechnika Gdańska, 2022.\newline\url{https://github.com/simle-stardust/Manual/blob/main/papers/2022_Jasiukowicz_praca_inzynierska.pdf}

\bibitem{blog2023} M.~Jasiukowicz, \emph{Stardust balloon mission from Iława}, 2023.\newline\url{https://yasiu.pl/blog/stardust-balloon-mission-from-ilawa/}

\bibitem{blog2022} M.~Jasiukowicz, \emph{SimLE Stardust Mission} -- wideo z lotu z 27.02.2022 i ślady trasy (KML), 2022.\newline\url{https://yasiu.pl/blog/simle-stardust-mission/}\newline Film: \url{https://www.youtube.com/watch?v=vfBzfLm3els}

\bibitem{go2023} Gazeta Olsztyńska, \emph{Absolwenci „Żeromka” puszczają balon (naukowy!)}, czerwiec 2023.\newline\href{https://gazetaolsztynska.pl/ilawa/940869,ILAWA-Absolwenci-quotZeromkaquot-puszczaja-balon-naukowy.html}{\nolinkurl{https://gazetaolsztynska.pl/ilawa/940869}}

\bibitem{infoilawa2023} infoIława, \emph{Zagadka rozwiązana. Tajemniczy balon prowadził nad Iławą eksperymenty}, maj 2023.\newline\url{https://www.infoilawa.pl/drukujpdf/artykul/71020}

\bibitem{storm} Most Wiedzy (Politechnika Gdańska), projekt \emph{STORM -- Stratospheric Toxin Migration}.\newline\url{https://mostwiedzy.pl/en/project/storm-migracja-toksyn-w-stratosferze,1159-1}

\bibitem{sed} Zespół Stardust, \emph{BEXUS 30 Stardust -- Student Experiment Documentation v5.0}, REXUS/BEXUS, 2022.\newline\url{https://rexusbexus.net/wp-content/uploads/2023/01/BX30_STARDUST_SED_v5-0s_26Apr2022_highlighted.pdf}

\bibitem{blogbexus} M.~Jasiukowicz, \emph{Jak dostaliśmy się do REXUS/BEXUS}, 2019.\newline\url{https://yasiu.pl/blog/jak-dostalismy-sie-do-rexus-bexus/}

\bibitem{bexus34} Polish Science, \emph{Students from Gdansk University of Technology qualify for the 34th edition of the REXUS/BEXUS programme}, styczeń 2024.\newline\url{https://www.polishscience.pl/en/?p=329276}

\bibitem{pg2021} Politechnika Gdańska, \emph{Students are sending a balloon into the stratosphere to study mysterious bacteria}, wrzesień 2021.\newline\url{https://pg.edu.pl/en/news/2021-09/students-are-sending-balloon-stratosphere-study-mysterious-bacteria}

\bibitem{bexusmanual} EuroLaunch (SSC, DLR MORABA), \emph{BEXUS User Manual}, v8, 2021.\newline\url{https://rexusbexus.net/bexus/bexus-user-manual}

\bibitem{magrian2018} S.~Magrian, K.~Pelzner, \emph{Układ pobierający próbki mikrobiologiczne w warunkach stratosferycznych}, projekt dyplomowy inżynierski, Wydział Elektroniki, Telekomunikacji i Informatyki, Politechnika Gdańska, opiekun: dr inż. P.~Fiertek, 2018.\newline\url{https://github.com/simle-stardust/Manual/blob/main/papers/2018_Magrian_Pelzner_praca_inzynierska.pdf}

\bibitem{stardustweb} Zespół Stardust, \emph{Studying microbiological life in the stratosphere} -- strona projektu z historią misji.\newline\url{https://stardust-project.netlify.app}

\bibitem{dok2016} Zespół Stardust, \emph{Stardust -- Balon stratosferyczny}, dokumentacja projektu przed pierwszym lotem, Koło Naukowe SimLE, Politechnika Gdańska, 2016.\newline\url{https://github.com/simle-stardust/Manual/blob/main/papers/2016_Stardust_dokumentacja_projektu.pdf}

\bibitem{gbsc2018} Zespół Stardust, zgłoszenie misji z 12.05.2018 do konkursu \emph{Global Balloon Space Challenge} (opis misji, dane z rejestratora, zdjęcia), 2018. Archiwum zespołu.
\bibitem{film2018} Zespół Stardust, film z misji z 12.05.2018, YouTube.\newline\url{https://www.youtube.com/watch?v=Gm_PSd67600}
\bibitem{film2022} Zespół Stardust NEXT, film z lotu szkoleniowego z 27.02.2022, YouTube.\newline\url{https://www.youtube.com/watch?v=vfBzfLm3els}
\bibitem{rs41nfw} \emph{rs41-nfw} -- przykład własnego oprogramowania dla radiosondy RS41.\newline\url{https://github.com/Nevvman18/rs41-nfw}
\bibitem{fbsimle} SimLE, posty o misjach Stardust w serwisie Facebook (2015--2018) i profil projektu Antares. Lista postów: \texttt{papers/facebook.md}.\newline\url{https://www.facebook.com/SimLEpl}, \url{https://www.facebook.com/AntaresBexusProject}
\bibitem{fbstardust} Stardust Microbiology Project, posty o misjach w serwisie Facebook (2020). Lista postów: \texttt{papers/facebook.md}.\newline\url{https://www.facebook.com/stardust.microbiology.project}

\bibitem{nasahb} S.~A.~Roth, E.~M.~Reed, C.~D.~Yoder, E.~A.~Chaffee, J.~Romualdez, \emph{Scientific Ballooning Handbook -- 50th Anniversary Edition}, NASA/TP-20240004989, NASA Goddard Space Flight Center, 2025.\newline\url{https://ntrs.nasa.gov/citations/20240004989}

\bibitem{bexustech} DLR, \emph{BEXUS Technical Overview} -- warunki lotu i wymagania dla eksperymentów (ciśnienie około \SI{10}{mbar}, praca bez chłodzenia konwekcyjnego).\newline\url{https://www.dlr.de/de/ar/downloads/dokumente/rw/BEXUS_Technical_Overview.pdf}

\bibitem{chaosvac} Zespół CHAOS (BEXUS, Uniwersytet w Kilonii), \emph{Vacuum Tests} -- test szczelności pojemnika ciśnieniowego w komorze próżniowej.\newline\url{https://www.bexus.org/?p=626}

\bibitem{ripstopdiy} Ripstop.pl, \emph{Spadochron DIY} -- szycie spadochronów półkulistych o średnicy \SI{1}{m} i \SI{4}{m} (autor: J.~Szostek, PWr in Space).\newline\url{https://ripstop.pl/spadochron-diy/}

\bibitem{chutemaker} rckTom, \emph{ChuteMaker} -- generator wykrojów klinów spadochronu (GPL-3.0).\newline\url{https://github.com/rckTom/ChuteMaker}

\bibitem{cansatgores} Club Robótica Granada, \emph{Democratising CANSAT II -- Parachute gores} (CC BY-SA 4.0).\newline\url{https://clubroboticagranada.github.io/DemocratisingCANSAT2/Parachute/gores}

\bibitem{bryceparachute} S.~Bryce, \emph{Spherical Parachute Pattern Generator}.\newline\url{https://www.scottbryce.com/parachute/spherical_parachute.html}

\bibitem{sp0lnd} \znak{SP0LND}.\newline\url{https://sp0lnd.pl/}

\bibitem{aghss} AGH Space Systems -- studenckie koło naukowe AGH.\newline\url{https://spacesystems.agh.edu.pl/}

\bibitem{nakkaxchute} R.~Nakka, \emph{1 metre Cross Parachute} -- budowa spadochronu krzyżowego, opis i dokument PDF z wymiarami.\newline\url{https://www.nakka-rocketry.net/xchute1.html}\newline\url{https://www.nakka-rocketry.net/articles/1mxchute.pdf}

\bibitem{direwolf} J.~Langner (WB2OSZ), \emph{Dire Wolf} -- programowy modem i TNC APRS, dokumentacja w katalogu \texttt{doc} (m.in. \emph{Raspberry Pi SDR IGate}).\newline\url{https://github.com/wb2osz/direwolf}

\bibitem{rtlsdrigate} RTL-SDR.com, \emph{Setting up a Raspberry Pi based APRS RX IGate with an RTL-SDR}.\newline\url{https://www.rtl-sdr.com/setting-up-a-raspberry-pi-based-aprs-rx-igate-with-an-rtl-sdr/}

\bibitem{vpdigi} \znak{SQ8VPS}, \emph{VP-Digi} -- digipeater APRS na STM32 (GPL-3.0, dokumentacja po polsku).\newline\url{https://github.com/sq8vps/vp-digi}

\bibitem{satnogsbuild} Libre Space Foundation, \emph{SatNOGS Wiki -- Build} (CC BY-SA).\newline\url{https://wiki.satnogs.org/Build}

\bibitem{klofasmobile} B.~Klofas, \emph{Mobile Horus Binary v2 and APRS Receiving Station with Raspberry Pi}, 2025 (CC BY 4.0).\newline\url{https://www.klofas.com/blog/2025/mobile-horus-aprs-receive-station-with-raspberry-pi/}

\bibitem{laaces} LaSPACE (Louisiana Space Grant Consortium), \emph{LaACES Student Ballooning Course} -- wykłady, ćwiczenia i szablony dokumentów (adres za~\cite{nasahb}).\newline\url{https://laspace.lsu.edu/laaces/student-balloon-course/}

\bibitem{nebp} Nationwide Eclipse Ballooning Project (Montana State University), kursy: \emph{Engineering Course} i \emph{Science Course} (adres za~\cite{nasahb}).\newline\url{https://eclipse.montana.edu/education/engineeringcourse.html}

\bibitem{sba} Stratospheric Ballooning Association, \emph{Academic High-Altitude Conference (AHAC)} i \emph{Journal of High Altitude Ballooning} (adres za~\cite{nasahb}).\newline\url{https://www.stratoballooning.org/ahac}

\bibitem{sera} Rozporządzenie wykonawcze Komisji (UE) nr 923/2012 (SERA), Dodatek 2: Bezzałogowe balony wolne.\newline\url{https://eur-lex.europa.eu/eli/reg_impl/2012/923/oj}

\bibitem{pazp} Polska Agencja Żeglugi Powietrznej, \emph{Rezerwacja przestrzeni powietrznej -- sondowania atmosfery, inne}.\newline\href{https://www.pansa.pl/asm1/rezerwacja-przestrzeni-powietrznej-sondowania-atmosfery-inne/}{\nolinkurl{https://www.pansa.pl/asm1/}}\newline\href{https://www.pansa.pl/asm1/rezerwacja-przestrzeni-powietrznej-sondowania-atmosfery-inne/}{\nolinkurl{rezerwacja-przestrzeni-powietrznej-sondowania-atmosfery-inne/}}\newline Formularz: \url{https://www.pansa.pl/wp-content/uploads/2025/07/Formularz_sonda_v1.pdf}

\bibitem{ksat} KSat e.V. (Kleinsatelliten), Uniwersytet w Stuttgarcie -- studencka organizacja kosmiczna, projekty balonowe BUBBLE i PARSEC.\newline\url{https://ksat-stuttgart.de/en/about-us/}

\bibitem{amu} Allianz Global Corporate \& Specialty, \emph{Allgemeine Versicherungsbedingungen für die Haftpflichtversicherung von Luftfahrzeugen (AMU 300/06)}.\newline\url{https://www.peschke-muc.de/wp-content/uploads/2012/08/AMU-300-06-Haftpflicht.pdf}

\bibitem{faqradiowe} Hackerspace Kraków, \emph{FAQ radiowe -- jak zostać krótkofalowcem}.\newline\url{https://github.com/HackerspaceKRK/faq-radiowe}

\bibitem{habduino} HABduino -- tracker dla balonów stratosferycznych.\newline\url{https://github.com/HABduino/HABduino}\newline Kopia instrukcji (CC BY-SA 3.0): \url{https://github.com/simle-stardust/Manual/tree/main/papers/habduino}

\bibitem{pits} Pi In The Sky -- tracker GPS do Raspberry Pi dla balonów stratosferycznych.\newline\url{https://www.pi-in-the-sky.com/}

\bibitem{flextrack} D.~Akerman, \emph{FlexTrack} -- oprogramowanie trackera (RTTY, LoRa, APRS).\newline\url{https://github.com/daveake/FlexTrack}

\bibitem{flextrackpr} \emph{Change default for APRS\_HIGH\_USE\_WIDE2} -- zmiana domyślnej ścieżki APRS w FlexTrack na NOPATH.\newline\url{https://github.com/daveake/FlexTrack/pull/15}

\bibitem{ukhasantena} UKHAS, \emph{Payload antenna} -- wykonanie anteny ćwierćfalowej z przeciwwagą.\newline\url{https://ukhas.org.uk/doku.php?id=guides:payload_antenna}\newline Kopia w Internet Archive: \url{https://web.archive.org/web/20210925035738/https://ukhas.org.uk/guides:payload_antenna}

\bibitem{akerman} D.~Akerman, \emph{Tracking a high altitude balloon}.\newline\url{https://www.daveakerman.com/?p=1732}

\bibitem{sondehubcalc} SondeHub, \emph{Burst Calculator} i \emph{Predictor}.\newline\url{https://sondehub.org/calc/}\newline\url{https://predict.sondehub.org/}

\bibitem{nasatypes} NASA, \emph{Scientific Balloons: Types of Balloons}, 2015.\newline\url{https://www.nasa.gov/scientific-balloons/types-of-balloons}

\bibitem{loon} X, the moonshot factory, \emph{Loon -- collection} oraz \emph{The Loon Library}.\newline\url{https://x.company/collection/loon/}\newline\url{https://storage.googleapis.com/x-prod.appspot.com/files/The%20Loon%20Library.pdf}

\bibitem{akademia} Akademia Stratosferycznie -- kursy o misjach stratosferycznych.\newline\url{https://akademia.stratosferycznie.pl/}

\bibitem{nist} NIST Chemistry WebBook, \emph{Thermophysical Properties of Fluid Systems} (hel: gęstość i współczynnik ściśliwości).\newline\url{https://webbook.nist.gov/chemistry/fluid/}

\bibitem{esabx34} ESA, \emph{BEXUS 34/35 Stratospheric Balloon Campaign: A Milestone for University Students}, 2024.\newline\url{https://www.esa.int/Education/Rexus_Bexus/BEXUS_34_35_Stratospheric_Balloon_Campaign_A_Milestone_for_University_Students}

\bibitem{nasahb1975} National Center for Atmospheric Research, \emph{Scientific Ballooning Handbook}, 1975 (NTIS PB-246575).\newline\url{http://www.xlta.org/library/noaa/Scientific_Ballooning_Handbook_Part1.pdf}

\bibitem{titan1} \emph{Titan-1 -- Design and Mission Document}.\newline\url{https://bombasaro.org/files/Titan1_Design_Mission_Doc.pdf}

\bibitem{rober} M.~Rober, \emph{Egg Drop From Space}, film, YouTube.\newline\url{https://www.youtube.com/watch?v=BYVZh5kqaFg}

\bibitem{stratocat} StratoCat -- archiwum lotów balonów stratosferycznych i biuletyn \emph{World Balloon News}.\newline\url{https://stratocat.com.ar}\newline\url{https://stratocat.substack.com}

\bibitem{ki4mcw} KI4MCW, \emph{Cutdown Mechanisms}.\newline\url{https://sites.google.com/site/ki4mcw/Home/cutdown-mechanisms}

\bibitem{r2home} Y.~Hadji, \emph{R2Home -- Making Radiosondes Reusable}.\newline\url{https://www.r2ho.me/}

\bibitem{supercluster} SuperCluster -- misje balonów stratosferycznych.\newline\url{https://www.supercluster.it/}

\bibitem{loraaprs} \emph{LoRa APRS Polska -- przewodnik po podstawach i standardach}.\newline\url{https://lora-aprs.pl/}

\bibitem{tbtracker} R.~Kroes, \emph{TBTracker-RX} -- przenośny odbiornik LoRa dla balonów.\newline\url{https://github.com/RoelKroes/TBTracker-RX}

\bibitem{m20fw} \znak{SQ2IPS}, \emph{m20-custom-firmware} -- własne oprogramowanie radiosondy M20.\newline\url{https://github.com/sq2ips/m20-custom-firmware}

\bibitem{wenet} Project Horus, \emph{Wenet} -- transmisja zdjęć z balonu.\newline\url{https://github.com/projecthorus/wenet}\newline Odbiór w przeglądarce: \url{https://horus.sondehub.org/}\newline Składanie obrazów: \url{https://ssdv.habhub.org/}

\bibitem{sp2gryf} Kluby \znak{SP2ZIE} i \znak{SP2YWL}, opis misji balonowej \znak{SP2GRYF} z 11.07.2026.\newline\url{https://sp2gryf.pl/balon/}

\bibitem{tinygs} TinyGS -- otwarta sieć stacji naziemnych LoRa.\newline\url{https://tinygs.com/}\newline\url{https://github.com/G4lile0/tinyGS}

\bibitem{meshtastic} Meshtastic -- sieć mesh na modułach LoRa.\newline\url{https://meshtastic.org/}

\bibitem{ogn} Open Glider Network -- sieć śledzenia szybowców i tracker ESP32.\newline\url{http://wiki.glidernet.org/}\newline\url{http://wiki.glidernet.org/esp32-ogn-tracker}

\bibitem{windborne} WindBorne Systems -- sieć autonomicznych balonów meteorologicznych z powłoką z folii.\newline\url{https://windbornesystems.com/}

\bibitem{ieee2025} J.~Dean, \emph{Inside the Best Weather-Forecasting AI in the World}, IEEE Spectrum, 21.10.2025. Zdjęcie: Christie Hemm Klok.\newline\url{https://spectrum.ieee.org/ai-weather-forecasting}

\end{thebibliography}
